\section{整数下段：素数证书、CHL 十二行与指数拼接}
\label{sec:lower}

解析方法覆盖的是以实变量 $x$ 表示的巨大区间；要得到所有正整数 $n$
的结论，还必须填满 $x=n^6$ 的较小范围。本节在唯一外部输入 H7 下证明
\begin{equation}\label{lower:target}
 1\leq n\leq N_*:=1242317999994
 \quad\Longrightarrow\quad
 \exists p\text{ 为素数},\quad n^6<p<n^6+6n^5.
\end{equation}
其中 $1\leq n\leq1259$ 完全由本库的实际素性证明处理；
其余范围使用 H7 中明确列出的十二条 CHL 短区间结论。
本节不使用 QRH、CS、Rosser 或有限 RH。

这一区分对理解证明的可信基础很重要。历史数值文件曾记录一种有限基素性测试；
当前 Lean 证明的这些素数已另由递归 Lucas 证书证明，因此该测试及其适用范围
并不是当前定理的额外文献输入。另一方面，CHL 行仍是 H7，
不能因为它们的整数转换已经形式化，就声称 CHL 的上游解析定理也已形式化。

\subsection{1259 个小区间的确定素数见证}

对每个整数 $1\leq n\leq1259$，保存一个具体自然数 $p_n$。
形式化需要分别证明三件事：
\begin{equation}\label{lower:witness-obligations}
 \operatorname{Prime}(p_n),\qquad n^6<p_n,\qquad
 p_n<n^6+6n^5.
\end{equation}
后两项是整数算术命题，不能由“$p_n$ 大致在 $n^6$ 附近”代替。
前几项及末项如下，完整列表由随附数据与源文件提供：
\begin{center}
\begin{tabular}{r|r|r}
$n$ & $n^6$ & $p_n$\\\hline
1 & 1 & 2\\
2 & 64 & 67\\
3 & 729 & 733\\
4 & 4096 & 4099\\
5 & 15625 & 15629\\
1259 & 3982487126873086441 & 3982487126873086493
\end{tabular}
\end{center}
完整见证列表在 \texttt{results/lower\_range\_k6\_certificate.json} 中；
与其对应的定理是 \texttt{PrimeWitnesses.lean} 中的
\texttt{small\_prime\_1} 至 \texttt{small\_prime\_1259}。
每条定理引用一个已经证明的素性命题，并直接证明两个严格整数不等式。

\begin{lemma}[本项目使用的 Lucas 素性准则]\label{lower:lucas}
设整数 $P\geq2$，并已知 $P-1$ 的完整素因数分解。
若存在模 $P$ 的元素 $a$，满足
\begin{equation}\label{lower:lucas-congruences}
 a^{P-1}\equiv1\pmod P,
 \qquad
 a^{(P-1)/q}\not\equiv1\pmod P
 \quad\text{对每个素因子 }q\mid P-1,
\end{equation}
则 $P$ 为素数。
\end{lemma}

\begin{proof}
第一个同余保证 $a$ 是模 $P$ 的可逆元，其阶整除 $P-1$。
如果该阶是 $P-1$ 的真因子，则 $(P-1)/\operatorname{ord}(a)$
有某个素因子 $q$，从而该阶整除 $(P-1)/q$，与第二个同余条件矛盾。
因此 $a$ 的阶恰为 $P-1$。由有限单位群的阶关系，
$P-1\leq\varphi(P)$；而 $P\geq2$ 时 $\varphi(P)\leq P-1$，
等号只可能在 $P$ 为素数时成立。
\end{proof}

实际实现使用 Mathlib 的 \texttt{lucas\_primality}。
给出 $P-1=\prod_iq_i^{e_i}$ 时，既检查乘法恒等式，
也递归引用每个 $q_i$ 的已证明素性；然后对所有这些素因子检查
式~\eqref{lower:lucas-congruences}。由于 $q_i<P$，依赖关系严格向较小整数
下降，形成可复用节点的有向无环证书。完整数据共有
\emph{5349 个素数节点}，分为 27 个源文件
\texttt{Pratt/Block000.lean} 至 \texttt{Pratt/Block026.lean}。
小于 $1000$ 的节点由 \texttt{norm\_num} 的有限算术素性证明处理，
其余节点用上述 Lucas 准则。

例如，对于见证 $P=15629$，采用
\[
 P-1=2^2\cdot3907,\qquad a=2.
\]
在已经证明 $2$ 与 $3907$ 为素数之后，需检查
\[
 2^{15628}\equiv1,\qquad
 2^{7814}\equiv15628\ne1,\qquad
 2^4\equiv16\ne1\pmod{15629}.
\]
这条证书可在 \leanref{Pratt/Block005}{prime_15629} 中直接核查。
其它较大的见证遵循同一逻辑，不要求读者相信一个外部程序输出的
\texttt{is\_prime=true}。

Python/FLINT 负责寻找因数分解和模幂见证并生成候选 Lean 源码。
这些候选只有在整数乘法、模幂和递归素性证明均通过内核检查后才能使用。
模幂计算由 \texttt{reduce\_mod\_char} 生成证明项；这里没有采用
\texttt{native\_decide} 作为捷径，也没有把 Python 的返回值注册成数学公理。

\begin{proposition}[小整数范围]\label{lower:small}
对每个 $1\leq n\leq1259$，存在素数 $p$ 满足
$n^6<p<n^6+6n^5$。
\end{proposition}

\begin{proof}
对有限个整数 $n$ 分情形，每个叶节点使用已经证明的
式~\eqref{lower:witness-obligations}。
实际源文件用二分情形树组织这些叶节点，以控制证明深度；
二分组织方式不改变被穷尽的整数集合。
\end{proof}

对应定理为 \leanref{PrimeWitnesses}{small_range_prime}。
特别地，$n=1$ 由 $p=2$ 处理，故后面要求 $n\geq2$ 的合数端点论证
不会遗漏这个最小正整数。

\subsection{CHL 的区间方向与严格端点转换}

H7 的每一行采用下列准确的方向和端点约定：给定 $\Delta>1$ 与 $X$，
\begin{equation}\label{lower:CHL-input}
 y\geq X\quad\Longrightarrow\quad
 \exists p\text{ 为素数},\quad
 y\left(1-\frac1\Delta\right)<p\leq y.
\end{equation}
源文为 \cite{CHL,CHLcorr} 的修订结论及指定表行。
这是左开右闭的区间；不能直接将右端点的不严格不等式抄成严格不等式。
接口 \leanref{Definitions}{CHLInterval} 保留了这一差别。

\begin{lemma}[单行 CHL 到六次幂短区间]\label{lower:CHL-conversion}
假设式~\eqref{lower:CHL-input} 成立。若整数 $n\geq2$ 满足
\begin{equation}\label{lower:row-conditions}
 X\leq n^6,\qquad n\leq6(\Delta-1),
\end{equation}
则存在素数 $p$，使 $n^6<p<n^6+6n^5$。
\end{lemma}

\begin{proof}
令 $y=n^6+6n^5=n^5(n+6)$。
因 $y\geq n^6\geq X$，可对 $y$ 使用 H7 的对应行。
关键的左端点比较为
\begin{align}
 y\left(1-\frac1\Delta\right)-n^6
 &=\frac{n^5}{\Delta}\bigl(6(\Delta-1)-n\bigr)\geq0.
 \label{lower:left-comparison}
\end{align}
于是源文给出的素数满足
$p>y(1-1/\Delta)\geq n^6$。
右端仅有 $p\leq y$，但 $n\geq2$ 时两个整数因子
$n^5$ 和 $n+6$ 都严格大于 $1$，所以 $y=n^5(n+6)$ 是合数。
由于 $p$ 是素数，$p\ne y$，因此 $p<y$。
\end{proof}

这两个步骤分别对应
\leanref{CHL}{nat_sixth_short_of_CHL} 与
\leanref{Reduction}{sixth_right_endpoint_not_prime}。
注意式~\eqref{lower:row-conditions} 容许等号：即使
$n=6(\Delta-1)$，左端仍由 CHL 的严格不等式保证 $p>n^6$。

\subsection{全部十二行及整数阈值}

下表列出当前 \leanref{CHLRange}{CHLRows} 的全部输入行，
并给出本项目使用的整数覆盖区间 $[a_j,b_j]$。
这里 $a_j$ 是已核验的安全下阈值，不声称是最小可能阈值；
$b_j=6(\Delta_j-1)$ 是精确整数。

\begin{center}
\small
\begin{tabular}{c|r|r|r}
$X_j$ & $\Delta_j$ & $a_j$ & $b_j=6(\Delta_j-1)$\\\hline
$4\cdot10^{18}$ & $39097000$ & $1260$ & $234581994$\\
$\e^{43}$ & $41816800$ & $1296$ & $250900794$\\
$\e^{46}$ & $163940000$ & $2136$ & $983639994$\\
$\e^{50}$ & $1061200000$ & $4161$ & $6367199994$\\
$\e^{55}$ & $10288400000$ & $9573$ & $61730399994$\\
$\e^{60}$ & $76918400000$ & $22027$ & $461510399994$\\
$\e^{75}$ & $174043000000$ & $268338$ & $1044257999994$\\
$\e^{90}$ & $184304000000$ & $3269018$ & $1105823999994$\\
$\e^{105}$ & $191886000000$ & $39824785$ & $1151315999994$\\
$\e^{120}$ & $197917000000$ & $485165196$ & $1187501999994$\\
$\e^{135}$ & $202553000000$ & $5910522064$ & $1215317999994$\\
$\e^{150}$ & $207053000000$ & $72004899338$ & $1242317999994$
\end{tabular}
\end{center}

每行使用两类不同的证据。$\Delta_j$ 和源域 $X_j$ 是 H7 的文献输入；
而将其转换为 $[a_j,b_j]$ 所需的
\begin{equation}\label{lower:threshold-checks}
 X_j\leq a_j^6,\qquad b_j\leq6(\Delta_j-1)
\end{equation}
均在本库中证明，未再假设一个整数区间结论。
第一行的下界是直接的整数计算：
\[
 1260^6=4001504141376000000>4\cdot10^{18}.
\]
其余十一条指数不等式完整列为
\begin{alignat}{2}
 \e^{43}&\leq1296^6,&\qquad \e^{46}&\leq2136^6,\notag\\
 \e^{50}&\leq4161^6,& \e^{55}&\leq9573^6,\notag\\
 \e^{60}&\leq22027^6,& \e^{75}&\leq268338^6,\notag\\
 \e^{90}&\leq3269018^6,& \e^{105}&\leq39824785^6,\notag\\
 \e^{120}&\leq485165196^6,&
 \e^{135}&\leq5910522064^6,\notag\\
 \e^{150}&\leq72004899338^6.&&
 \label{lower:eleven-exp-thresholds}
\end{alignat}
这些比较由 \texttt{Thresholds.lean} 中
\texttt{chl\_threshold\_43} 至 \texttt{chl\_threshold\_150}
的反射式二进区间证明验证，设置的工作精度为 40 位。
指数的数值包围本身生成内核可检查的证明，
并非先在外部取 $\lceil\e^{j/6}\rceil$ 再把这个整数当作正确答案。

如果 $a_j\leq n\leq b_j$，幂函数在非负实轴的单调性给出
$X_j\leq a_j^6\leq n^6$，同时 $n\leq6(\Delta_j-1)$。
因此引理~\ref{lower:CHL-conversion} 对该行覆盖的每个整数都适用。
形式化把这个共同步骤封装成
\leanref{CHLRange}{nat_short_of_CHL_integer_row}。

\begin{proposition}[CHL 行之间没有整数空隙]\label{lower:CHL-coverage}
在 H7 下，对每个整数
$1260\leq n\leq1242317999994$，存在素数 $p$ 满足
$n^6<p<n^6+6n^5$。
\end{proposition}

\begin{proof}
按表中顺序记十二个区间为 $[a_j,b_j]$，$0\leq j\leq11$。
表内精确整数满足
\[
 a_0=1260,\qquad b_0<b_1<\cdots<b_{11}=N_*,
 \qquad a_{j+1}\leq b_j+1\quad(0\leq j<11).
\]
所以 $[a_0,b_0]$ 与其后逐行增加的区间覆盖整个整数区间
$[1260,N_*]$。等价地，对给定的 $n$，取第一个满足 $n\leq b_j$ 的行；
若 $j>0$，前一行未覆盖便给出 $n\geq b_{j-1}+1\geq a_j$，
从而这一行的上下条件都成立。
每个比较都可从表中整数直接验证；无需关于素数分布的进一步输入。
\end{proof}

对应定理为 \leanref{CHLRange}{chl_integer_coverage}。
与命题~\ref{lower:small} 拼接即得
\leanref{LowerRange}{lower_range_complete}，也就是式~\eqref{lower:target}。
这里的 $1259/1260$ 接口是连续的；我们没有以“有限多个剩余例外可计算”
结束论证，而是给出了这些例外的实际素数证明。

\subsection{与 \texorpdfstring{$L=167$}{L=167} 的精确拼接}

要与后面的实变量中段相接，还需证明式~\eqref{lower:target} 至少覆盖
所有满足 $\log(n^6)\leq167$ 的正整数。这依赖一条另外验证的超越不等式：
\begin{equation}\label{lower:join-exp}
 \e^{167}<(1242317999994)^6=N_*^6.
\end{equation}
\leanref{Thresholds}{lower_join_exp} 使用 20 位反射式二进区间证明建立该式。
它与上一组十一条阈值检查一样，是 Lean 中的实数不等式证明。

\begin{lemma}[指数接口]\label{lower:log-join}
设 $n\in\N$ 且 $n\geq1$。若 $\log((n:\R)^6)\leq167$，则 $n\leq N_*$。
\end{lemma}

\begin{proof}
由 $n>0$，指数函数的单调性给出
\[
 n^6=\exp(\log(n^6))\leq\e^{167}<N_*^6.
\]
若 $n>N_*$，非负整数的六次幂保持序关系，将与最后一个严格不等式矛盾。
所以所需的非严格上界 $n\leq N_*$ 成立。
\end{proof}

\begin{theorem}[完整整数下段]\label{lower:log-theorem}
在 H7 下，对每个 $n\in\N$，若 $n\geq1$ 且
$\log((n:\R)^6)\leq167$，则存在素数 $p$ 满足
\[
 n^6<p<n^6+6n^5.
\]
\end{theorem}

\begin{proof}
先由引理~\ref{lower:log-join} 得 $n\leq N_*$，再使用
小范围证书与 CHL 行已经拼成的式~\eqref{lower:target}。
\end{proof}

对应实现为 \leanref{CHLRange}{lower_join_of_log} 和
\leanref{LowerRange}{lower_log_range}。
式~\eqref{lower:join-exp} 为严格不等式，说明两个范围有重叠；
并不需要猜测 $\e^{167/6}$ 位于哪两个相邻整数之间。

\subsection{H7 的来源状态与本节证明的边界}

CHL 原文与勘误均有正式期刊发表记录。当前十二行与作者 arXiv v3 修订本的
数值、定义域以及左右端点约定相符；但是，这与独立验证其全部上游计算不是同一件事。
源文复核发现，公开的有限位辅助参数在
$\log X=50,60,75,90,105,135$ 六行生成的 $\Delta$
略小于表中目标值。该方向差异已经用以印刷十进制数为有理输入的 Arb 计算确认。
由于较大的 $\Delta$ 对应较短的区间，不能用略小的生成值直接证明表中较强的行。
小幅调整辅助参数后能够达到目标的数值观察，也不代替完整严格认证。

本节因此明确接受的是 H7 中的十二条\emph{区间命题}，
不是声称已由公开舍入见证在本库重证这些命题。
本节对它们之后的步骤——整数阈值、各行覆盖、素数证书、严格端点和指数拼接——
给出了完整证明。正式勘误正文的精确表格比对以及 CHL 判据的完整独立认证仍列在来源审查中。
这一限制应与定理同时阅读，而不能被“所有下游模块编译通过”一句话覆盖。
